\clearpage
\section{Cross-phase statistical and validity plan}

\subsection{Confirmatory hierarchy}

The report should use sequential gates rather than treating every endpoint as independently confirmatory:
\begin{enumerate}
  \item frozen stage-aware ranking transfers to new parents;
  \item complete snapshot replay makes counterfactual branches valid;
  \item candidate/operator sets contain recovery headroom;
  \item a recovery-value critic ranks candidate utility on sealed data;
  \item critic selection works prospectively under a fixed trigger;
  \item the complete utility system improves final task success;
  \item external validity and flow guidance are tested only after the core system passes.
\end{enumerate}

A failed detector gate does not prohibit studying recoverability with time, progress, and observable context. A failed oracle-opportunity or critic-actionability gate does prohibit claiming a useful local recovery critic or gradient guide.

\subsection{Alarm and timing definitions}

For failed unit $u$, let $H_u$ be the frozen normalization horizon and $\tau_u$ the first alarm. Define miss-adjusted detection time
\begin{align}
T_u &= \min(\tau_u/H_u,1) && \text{if detected},\\
T_u &= 1 && \text{if missed}.
\end{align}
The primary timing contrast is $\mathrm{E}[T^{\mathrm{SAFE}}-T^{\mathrm{time}}]$; negative favors SAFE. Conditional-on-detection time is secondary because a method could appear early by detecting only a few easy failures.

For $q\in\{0.10,0.25,0.50\}$, early-landmark recall is
\begin{equation}
  \Pr(\tau_u/H_u\le q\mid u\text{ fails}).
\end{equation}
The successful-parent false-intervention rate is the probability of at least one alarm/intervention over an entire paired would-be successful parent. A treated rollout alone does not reveal whether it would have succeeded without intervention; this outcome requires a matched no-intervention clone or randomized arm under the same initial identity.

\subsection{Calibration}

For successful-parent maximum scores $Z_1,\ldots,Z_n$ and target $\alpha$, a split-conformal upper threshold uses order statistic
\begin{equation}
  k=\lceil(n+1)(1-\alpha)\rceil.
\end{equation}
If $k>n$, place the threshold above the observed maximum and label the operating point unresolved. Zero empirical FPR with too few calibration successes is not evidence of precise control.

Calibration requirements are:
\begin{itemize}
  \item successful parent is the calibration unit;
  \item score aggregation, tie handling, task weighting, threshold rule, and target $\alpha$ are frozen;
  \item normalization is fitted on training, not calibration or test;
  \item calibration and final evaluation identities are disjoint;
  \item report calibration count, order-statistic rank, realized FPR, and exact or Wilson interval;
  \item evaluate drift by task/layout/policy without updating the final-test threshold.
\end{itemize}

\subsection{Randomization and paired inference}

Intervention assignments should be generated before outcomes and blocked by task, layout/style, seed/reset family, and server/GPU session. Arms must be interleaved within server sessions. The primary analysis is intention to treat: differential runtime or model errors count against the assigned arm according to the frozen attrition rule.

For exactly paired initial identities, report benefit and harm discordance:
\begin{align}
  B &= \Pr(Y^{\mathrm{method}}=1,Y^{\mathrm{control}}=0),\\
  H &= \Pr(Y^{\mathrm{method}}=0,Y^{\mathrm{control}}=1),\\
  \Delta_{\mathrm{paired}} &= B-H.
\end{align}
Use blocked randomization/permutation inference as primary where assignment supports it, task-then-parent hierarchical bootstrap as a transparent interval, and mixed-effects logistic regression only as a sensitivity analysis.

\subsection{Bootstrap hierarchy}

For detector comparisons on identical parents:
\begin{enumerate}
  \item resample tasks if claiming task-population generalization;
  \item resample parents within each selected task;
  \item retain every stage, prefix, snapshot, and candidate from the sampled parent;
  \item recompute the complete metric and paired method contrast.
\end{enumerate}
For candidate studies, snapshots are nested within parents and candidate branches remain paired within snapshot. Three neural initialization seeds measure fitting instability and do not replace task or parent uncertainty.

\subsection{Power and sample size}

Final $N$ should be frozen through simulation or hierarchical bootstrap using development or blinded-pilot variance. Allowed adaptations include extending a fixed-attempt collection because predeclared support counts are unmet and blinded variance-based re-estimation. It is not acceptable to stop when ROC or success first becomes favorable, add tasks after seeing per-task effects, or alter thresholds from test alarms.

Planning calculations from current evidence suggest:
\begin{itemize}
  \item roughly 400 failures before clustering to detect a 0.05-horizon paired timing effect under the observed variance, and approximately 500-700 after heterogeneity inflation;
  \item roughly 330 parents per arm before clustering for a 12\%-to-20\% success difference, potentially 500-650 per main arm after task/server inflation;
  \item at least 200 successful calibration parents for a 5\% threshold, with approximately 400 preferable;
  \item approximately 25 successes and 25 failures per task-stage stratum for exploratory stage estimates, and closer to 50 per class for stable confirmatory estimates.
\end{itemize}

These calculations are not promises or venue rules. The final powered design should prioritize the primary two-arm comparison in Phase 7; secondary trigger and recovery ablations can use smaller frozen cohorts.

\subsection{Leakage and survivor-bias controls}

\begin{longtable}{P{0.24\textwidth}P{0.68\textwidth}}
\caption{Cross-phase validity controls.}\label{tab:validity-controls}\\
\toprule
Risk & Required control \\
\midrule
\endfirsthead
\toprule
Risk & Required control \\
\midrule
\endhead
Parent leakage & Assign parents before stages, windows, snapshots, or candidate augmentation. Every descendant inherits the parent split. \\
Snapshot leakage & Keep all operators/candidates from a snapshot together and weight snapshots, not candidate rows. \\
Future information & Exclude final horizon, termination reason, future predicate states, future stage completion, future features, and realized future duration from deployed inputs. \\
Privileged stage & Report simulator predicate stage and causal observable stage tracker separately. \\
Survivor bias & Publish risk-set denominator and class support at every later prefix; treat stage completion and failure as competing events in sensitivity analysis. \\
Outcome quota & Use case-control/outcome-quota data for ranking or feasibility only, never deployment prevalence, FPR burden, or success rate. \\
Feature duplication & Store one feature per genuine policy inference, not every cached action. \\
Restore confounding & Restore simulator, policy queues/cache, counters, instruction, tracker, and RNG; prove exact same-request/action replay. \\
Compute confounding & Match candidate count, recovery horizon, NFE, policy calls, or wall-clock budget according to the stated comparison and report both calls and time. \\
Oracle mixing & Never average oracle stage, oracle onset, environment rewind, or oracle candidate choice with deployable results. \\
Selective attrition & Freeze exclusion/retry rules; preserve partial manifests and technical errors; provide assigned/started/completed/invalid/analyzed flow. \\
\bottomrule
\end{longtable}

\subsection{Multiplicity and reporting}

Declare one primary contrast per phase and one final primary contrast: full utility system versus strongest frozen non-oracle baseline on final task success. Test secondary families in a frozen hierarchy or use Holm correction. Per-task, per-stage, per-delay, and failure-type analyses remain descriptive unless explicitly powered and adjusted. Report effect sizes and intervals even when gates fail.

\cautionnote{Gate interpretation.}{Pilot thresholds such as 20\% mixed snapshots, 15-point oracle headroom, pairwise ROC 0.65, or a 10-point critic gain are proposed continuation gates. They protect compute and scientific focus; they are not universal statistical truths.}

\subsection{Claim ledger}

\begin{longtable}{P{0.30\textwidth}P{0.23\textwidth}P{0.39\textwidth}}
\caption{Claim-to-evidence ledger at the current starting point.}\label{tab:claim-ledger}\\
\toprule
Claim & Current status & Evidence required \\
\midrule
\endfirsthead
\toprule
Claim & Current status & Evidence required \\
\midrule
\endhead
Xiaomi SAFE modestly improves matched-FPR classification. & Supported in frozen prospective shadow evaluation. & Preserve modest wording and artifact provenance. \\
Xiaomi SAFE provides reliable overall early warning. & Rejected by timing interval and early recall. & New representation plus untouched prospective timing cohort. \\
Full-parent stage conditioning improves early ranking. & Rejected as a general five-task claim: 0.696 development ROC fell to 0.489 in frozen identity-disjoint confirmation. & Preserve task/stage heterogeneity as diagnostic evidence; do not refit on confirmation. \\
Proprioceptive state context improves early stage ranking. & Unsupported: context reached 0.566 on the retrospective 85-parent holdout versus stage 0.584 and time 0.500. & A new untouched cohort would be required for a context claim, but it is not the current priority. \\
SAFE ranks useful action chunks. & Engineering path only. & Phase 3 opportunity and Phase 6 prospective selection. \\
Recovery-value critic predicts causal utility. & Planned. & Phase 5 sealed evaluation and Phase 6 prospective result. \\
Utility triggering improves task success. & Unsupported. & Phase 7 randomized closed-loop trial. \\
Flow guidance improves recovery. & Unsupported. & Phase 8 model-side, matched-compute execution. \\
Method generalizes to unseen tasks or policies. & Unsupported. & Phase 9 with identities unseen to critic development and separate policy calibration. \\
Method is real-robot deployable. & Unsupported; simulator predicates/resets are currently used. & Causal stage tracker and physical intervention study. \\
\bottomrule
\end{longtable}
